\section{Scope and reproducibility of the checks}
The algebraic reductions and physical branches in the preceding chapters
are fixed before any numerical comparison. The command \texttt{make check}
regenerates its required analytic intermediates and runs the standard
verification suite. It checks source consistency, exact coefficient
identities, amplitudes, integrated sectors, endpoint matching, asymptotic
expansions and numerical implementations. Each check records its exit
status and its output in \nolinkurl{build/checks/}. The command
\texttt{make check-list} lists the preparation steps and tests.
The more expensive full phase-space audits have separate commands in
the build instructions. Their saved results below are distinguished
from the standard suite.

For example, \nolinkurl{compton_dilog_identities.py} establishes the two
real dilogarithm reductions by differentiating both sides and fixing
their common boundary values at $s=1$. Only after that proof does
\nolinkurl{check_compton_dilog_identities.py} compare their implementations
at $s=1.0001,1.02,1.6,10,1000$, using 70-digit arithmetic and an
absolute tolerance of $10^{-60}$. These samples determine neither the
identity nor its integration constant.

The regular product form has a separate implementation in
\nolinkurl{compton_total_regular.py}. Its exact checks reconstruct every
coefficient and verify the real-axis inversion relations by derivatives
and boundary values. The numerical program
\nolinkurl{check_compton_total_regular.py} then compares the full result
with the original collected expression at ten energies from $X=0.0002$
to $X=200000$, including $X=1$ and offsets of $10^{-20}$ on either side.
Its tolerance is $10^{-30}\max(1,|C_1|)$. The paired trilogarithm is also
checked against its original boundary-value representation at those
points. These tests check the implementation of already established
identities.

Near threshold, subtracting primitive endpoints and dividing by $X^3$
amplifies rounding errors. The regular evaluator adds six working digits
per decade of $X<1$ to its base guard precision, then repeats the
evaluation with thirty additional digits. It checks both the change
and the imaginary residual relative to the result. A further increase
is allowed, and failure to converge raises an error.
At $X=10^{-18}$ and $10^{-24}$, the same test compares the result
with the independently derived low-energy expansion through $X^4$.
The relative differences are below $7\times10^{-55}$ and
$7\times10^{-73}$, respectively. These tests validate the evaluation
of the exact expression, without determining its coefficients.

The same program checks the shared root-function evaluator in
\nolinkurl{compton_root_classical.py}. It evaluates the common endpoint
differences once and uses the shuffle identity for the outer-zero word.
The original primitives, endpoint constants and complex branches are
unchanged. At the same ten energies, seven words are compared with their
separate implementations using 80-digit arithmetic and tolerance
$10^{-70}\max(1,|G|)$. The word $G(r,r,-1;X)$ is checked through its
defining derivative, with tolerance $10^{-65}\max(1,|\partial_XG|)$.

The direction averages have a separate check in
\nolinkurl{compton_real_analytic_angular_check.py}. It compares the closed
masters with direct sphere quadrature at four regular kinematic points.
Near threshold, small Gram determinants and cancellation between terms
require additional working precision. At $H=0.02$, the check repeats the
master sum at increasing precision and the sphere integral at two angular
resolutions. Both quadratures agree with the master sum to a relative
tolerance of $10^{-30}$. These evaluations share the same 25 coefficients,
so the comparison tests the direction averages and their implementation.
Precision and angular-resolution changes provide stability checks,
not rigorous error bounds or independent checks of the original trace.

\section{Independent amplitude checks}
The program \texttt{check\_compton\_amplitudes.py} constructs explicit
Dirac matrices with $m=e_{\mathrm T}=1$. It does not insert a known
squared Compton amplitude. It sums physical transverse photon
polarizations and performs the spin traces. Three Born points with
$(\omega/m,\cos\theta)=(0.1,-0.4),(1,0.3),(20,-0.7)$ agree with
Klein--Nishina to a relative error below $5\times10^{-16}$.

For the real amplitude it constructs on-shell three-body events
at $\omega/m=0.3,1,7$. At every event it replaces each of the three
photon polarizations in turn by its momentum. The spinor-sandwiched
matrix vanishes with a relative residual below $6\times10^{-14}$.
Interchanging the two outgoing photons changes the kernel by less
than $3\times10^{-16}$ relative to its norm. A separate test of the
contracted off-shell vertex integrand verifies
\eqref{eq:comp-supp-ward} before counterterms, with a relative
residual below $2\times10^{-16}$.

The soft check is more restrictive than a Ward identity. Keep a
hard-photon direction fixed and compare the exact three-body square
with $e_{\mathrm T}^2(-J^2)B_{\mathrm{amp}}$, projecting the hard system onto
on-shell two-body kinematics in the denominator. Table~\ref{tab:comp-supp-soft}
shows the approach to the derived limit. Corrections are linear in
the small photon energy at this fixed event.
\begin{table}[!htbp]
\centering
\begin{tabular}{cc}
\toprule
Soft photon energy divided by $m$ & Exact square / soft approximation\\
\midrule
$10^{-3}$ & $0.9980259445$\\
$10^{-4}$ & $0.9998019242$\\
$10^{-5}$ & $0.9999801857$\\
$10^{-6}$ & $0.9999980193$\\
\bottomrule
\end{tabular}
\caption{Soft-limit check of the six real-emission diagrams at
incident energy $\omega=m$. The event directions are fixed in the
verification script.}
\label{tab:comp-supp-soft}
\end{table}
The program \texttt{check\_compton\_virtual.py} separately integrates
the full subtracted loop tensor by Gauss quadrature. It uses
$\omega/m=0.3$, $\cos\theta=0.2$ and $\lambda/m=0.7$ so that
the parameter denominators stay away from physical cuts. This
finite regulator is a test of the matrix identities, not a physical
small-$\lambda$ prediction. Including all eight diagrams, the two
external Ward residuals decrease from $4.1\times10^{-9}$ at
quadrature order 12 to $7.6\times10^{-15}$ at order 28. The loop
matrix changes by $5.3\times10^{-15}$ relative to its norm between
orders 20 and 28. This independently tests the parameter reduction,
finite counterterms and relative signs. It does not evaluate the
inclusive total or the physical small-$\lambda$ loop across its cuts.

\section{Cross-check of the analytic scalar box}
The script \texttt{check\_compton\_scalar\_analytic.py} evaluates
\eqref{eq:comp-supp-box-dilog} at 50 decimal digits and compares it
with quadrature of \eqref{eq:comp-supp-box-one-dimensional}.
At $-t=0.001,0.09,5,7800$, the absolute differences are below
$3\times10^{-50}$. Quadrature serves only to check the derived
dilogarithm identity. This test checks the scalar box identity, rather than the total
phase-space integration.

\section{Cross-check of the closed vertex}
The script \texttt{check\_compton\_vertex\_analytic.py} compares
the six closed moments in \eqref{eq:comp-supp-vertex-closed-moments}
with their parameter integrals at $S=-4,0,0.49,1.6,10$.
Both real and imaginary parts are included. At 60 decimal digits,
the largest absolute difference is below $10^{-58}$.
The continuation in the integrand and in the closed expression
is fixed by the Feynman denominator as derived above.

The same program compares all four vertex insertions in the
Compton amplitude with parameter quadrature at
$(s,c)=(1.6,0.2),(10,-0.4),(100,0.7)$.
Increasing the quadrature order from 32 to 64 and 128 decreases
the maximum matrix discrepancy at the last point from
$8.33\times10^{-9}$ to $1.37\times10^{-10}$ and
$2.19\times10^{-12}$. Here the discrepancy is the Frobenius norm
of the matrix difference divided by the larger of one and the
norm of the closed matrix. The analytic expression itself has
no quadrature order. These tests concern the vertex contribution,
rather than the complete total-cross-section integration.

\section{Analytic integration of a virtual subset}
The trace coefficients for the vertex and self-energy terms are
generated by \nolinkurl{compton_virtual_trace_analytic.py} using
Clifford anticommutation. Explicit covariant polarization sums of
the Dirac matrices reproduce those coefficients to relative errors
below $3\times10^{-15}$ at the three points used above.
The covariant sum is used on both sides of this check because this
subset alone is not gauge invariant.

The primitive in \eqref{eq:comp-vt-u-primitive} has an exact symbolic
derivative equal to the crossed-ordering trace. Independent
50-digit quadrature of that trace at $s=1.6,10,100$ differs from
the primitive's endpoint difference by less than $10^{-45}$.
The script \nolinkurl{check_compton_virtual_angular_analytic.py}
also verifies the complete vertex-plus-self-energy angular integral,
including its phase-space normalization, against a separate
double-precision trace quadrature. At $\lambda=10^{-6}$ the analytic
coefficients $\mathcal C_{\mathrm{vertex}+\Sigma}$ are
\begin{align}
 s=1.6 &: \quad 52.9138831571782534624971544028,\\
 s=10 &: \quad 15.1114502033743705973520745024,\\
 s=100 &: \quad 2.81971068757998578698139625584.
\end{align}
These regulator-dependent numbers are diagram-subset checks.
They are not physical NLO correction factors and must be combined
with the spanning and real terms before comparison with a total
cross section.

\section{Ward identities of the complete closed amplitude}
The script \nolinkurl{compton_amplitude_analytic.py} combines the
closed self-energy, vertex and spanning matrices. No loop-parameter
quadrature is used. Replacing either external polarization by its
photon momentum gives normalized Ward residuals below
$1.1\times10^{-13}$ at
$(s,c)=(1.6,0.2),(10,-0.4),(100,0.7)$, for both
$\lambda=10^{-5}$ and $10^{-8}$. The normalization is the larger
of one and the norm of a transverse spinor-sandwiched amplitude,
as in the earlier Ward check. This tests the relative signs and
finite terms of the complete analytic virtual amplitude in the
physical region. It is distinct from the real-radiation energy integrals in the
total cross section.

The script \nolinkurl{check_compton_infrared_analytic.py} separately
changes $\lambda$ from $10^{-5}$ to $10^{-8}$ in this same closed
amplitude. At the three points above it verifies
\begin{equation}
 \frac{\operatorname{Re}\mathcal T(10^{-5})
       -\operatorname{Re}\mathcal T(10^{-8})}{4\ln(10^3)}
       =\mathcal B(b_e-1)
\end{equation}
with absolute residuals below $2\times10^{-14}$. This is the
coefficient required for pointwise cancellation against the
massive-photon soft term. It checks the independently derived
infrared factor, rather than determining that factor by fitting.

\section{Checks of the integrated amplitudes}
The following checks test the analytic integrations derived in Chapter 3.
Numerical values do not fix their constants or physical branches.

\subsection{Checks of the spanning matrix}
Numerical integration of the earlier two-parameter formula provides
a check of the analytic box, not its definition. At
$(s,c)=(1.6,0.2),(10,-0.4),(100,0.7)$, comparison of all
four physical polarization matrices in both photon orderings with
128-point integration in each parameter gives maximal relative
matrix differences $1.9\times10^{-13}$,
$7.1\times10^{-13}$ and $5.0\times10^{-12}$, respectively.
Separately, evaluating the rational interference coefficients
agrees with explicit covariant polarization contractions of the
closed $4\times4$ matrices to relative precision better than
$7.1\times10^{-15}$ at these points. Neither comparison uses a
published total-cross-section formula.

\subsection{Evaluation of the spanning angular primitives}
The numerical evaluation of the GPL expression can be performed
by convergent Taylor expansions without integrating the original
angle-dependent correction. Around a point $v_0$ away from a
letter $a$, write the Taylor coefficients of the shorter word
as $t_n$ and those of the derivative of $G(a,w)$ as $d_n$.
Its differential equation gives the recursion
\begin{equation}
 d_{-1}=0,\qquad
 d_n=\frac{t_n-d_{n-1}}{v_0-a},\qquad
 [ (v-v_0)^{n+1}]G(a,w;v)=\frac{d_n}{n+1}.
\end{equation}
Overlapping disks along the upper path then continue the same
analytic function to the endpoint. This operation evaluates the closed special functions on the
branch already fixed by the Feynman pole displacement and the
continuation path derived above. No cross-section comparison
enters that choice.

As a check, take $\lambda=10^{-5}$ and define
$J_j=\mathcal F_j(v_{\max};x)-\mathcal L_j(x)$ for $j=s,u$.
This is the quantity before multiplication by $1/(2x^2)$. Evaluation of the closed expressions gives
\begin{align}
 J_s(0.6)&=-18.6105396630505172566260787241204817816433900769,\\
 J_u(0.6)&=-14.1391278522157817227769489975337484383516616217,\\
 J_u(9)&=-3080.09775949773657530083054013237461366392218650.
\end{align}
Independent quadrature of the already derived angular interference
agrees with these values to absolute differences
$9\times10^{-47}$, $10^{-62}$ and $2\times10^{-58}$,
respectively. The last GPL evaluation crosses the apparent
$\eta$ pole by the specified upper path. Its residual imaginary
part is below $3\times10^{-64}$ before taking any real part.
These checks test the analytic integrations and the causal branch
choice. They do not use a published total cross section.

\subsection{Checks of the finite soft function}
The script \texttt{compton\_soft\_analytic.py} evaluates
\eqref{eq:comp-soft-K-dilog} directly and separately integrates
\eqref{eq:comp-soft-I-definition} for a check. At 70 decimal digits,
checks at $s=1.6$, $10$, $100$, and $19570.511835591835$, with
$a^2/q^2=0.0001$, $0.2$, $0.8$, and $0.9999$, give absolute differences below
$1.2\times10^{-66}$. These quadratures check the derived closed expression.
They do not define it. The result completes the finite soft parameter
integral at fixed angle. Its Born-weighted elastic angular integral
is a further part of the total-cross-section calculation.

\subsection{Check of the integrated infrared coefficient}
Direct checks of the
original angular integral at 65-digit working precision give
$\mathcal C_{\mathrm{IR}}=0.2034909958483591$, $10.14476601141435$,
$64.23179863499513$, and $359.2271977981074$ at $s=1.6$, $10$, $100$, and $19570.511835591835$. The largest absolute difference between the dilogarithmic
expression and quadrature is $1.1\times10^{-61}$.

\subsection{Independent numerical angular integration}
The numerical angular check evaluates the covariant Dirac trace before
integrating over the scattering angle.  Its scalar one-loop functions are
the logarithms and dilogarithms obtained from the parameter integrals.
It does not evaluate an integrated GPL primitive inside the numerical
integrand. In units $m=1$, put $x=s-1$, $y=x+t$, and use the
spin-averaged interference trace $\mathcal T$ and the Born factor
$\mathcal B=\overline{|\mathcal M_0|^2}/e_{\mathrm T}^4$ from
\eqref{eq:comp-supp-independent-total}.
The finite virtual-plus-soft coefficient in the local real-emission
subtraction scheme is
\begin{equation}
 \mathcal C_{\mathrm{VS}}(s)=\frac{1}{s}\int_{-1}^{1}\dd c\,
 \left\{\frac{\operatorname{Re}\mathcal T(s,c;\lambda)}4
 +\mathcal B(s,c)\left[(b_e-1)\ln\frac2\lambda-K(s,c)\right]\right\}.
\end{equation}
The logarithm of the photon mass cancels inside this combined integrand.
This coefficient must be combined with the real contribution defined by
the same local subtraction.  The finite coefficient obtained with the
analytic factor $(xy)^\eta$ has a different soft matching term. Here
$x,y$ in the regulator are the real-emission Dalitz variables, not the
elastic invariants used in this angular check.

To resolve both angular endpoints at large $s$, write $L=\ln s$ and
introduce $0<u<1$ by
\begin{align}
 c&=\left(1+\frac2x\right)\tanh\bigl[L(u-\tfrac12)\bigr],
 &y&=x-\frac{x^2}{2s}(1-c),\\
 \frac{dc}{du}&=\left(1+\frac2x\right)
 \frac{L}{\cosh^2[L(u-\tfrac12)]}.
\end{align}
The identity $\tanh(L/2)=x/(x+2)$ proves that the endpoints map exactly
to $c=-1$ and $c=1$.  Positive Gauss--Legendre nodes lie strictly inside
this interval, so no removable endpoint expression is evaluated directly.
The check uses arbitrary-precision arithmetic and three successively
doubled node counts.  Ten times the larger of the final two differences
is an empirical discretization error estimate, not a rigorous error bound.
The reported uncertainty also has the conservative arithmetic floor
$10^{20-d}$ when using $d$ decimal digits.
Increasing the arithmetic precision provides a separate check of
cancellation error, particularly near threshold.  A second variable map,
$y=(x/s)\exp(u\ln s)$, checks the angular Jacobian independently.
The scope of this test is the angular integration and its implementation.
The earlier comparison with parameter quadrature tests the scalar
one-loop functions themselves.

\subsection{Checks of the soft angular primitive}
The algebraic reduction has also been implemented independently of the
quadrature check. At $s=10$ it generates 306 terms involving 305 distinct
nonempty polylogarithm words, including words needed to evaluate their
lower-weight derivatives. Partial-fraction reconstruction of the
rational kernels agrees to $1.3\times10^{-60}$ at 60-digit precision.
At $z=0.9Z$, evaluating the generated polylogarithms through their defining
differential equations gives $17.00560385716334$, compared with
$17.00560385716422$ from direct quadrature of the rationalized kernel.
These two numerical evaluations check the algebraically integrated expression.

\subsection{Series evaluation of the soft angular endpoint}
\paragraph{Arbitrary-precision polylogarithm evaluation.}
The generated expression can be checked beyond machine precision by
convergent power series, without integrating the Compton kernel. For a
word $(a,\boldsymbol b)$ expand $G(a,\boldsymbol b;z_0+h)$ in powers of
$h$ with coefficients $c_n$, and write the tail word as
$G(\boldsymbol b;z_0+h)=\sum_n d_nh^n$. Its differential equation gives
\begin{equation}
 c_{n+1}=\frac{d_n-nc_n}{(z_0-a)(n+1)},\qquad z_0\ne a.
\end{equation}
Words are processed in increasing length, so all $d_n$ are already known.
At the regular origin of a word beginning with zero, $d_0=0$ and
$c_n=d_n/n$ for $n\geq1$. Steps smaller than the distance to the nearest
letter analytically continue these series along the real physical path.
At 60-digit working precision, the 306-term expression at $s=10$ and
$z=0.9Z$ agrees with direct quadrature to $1.0\times10^{-59}$.
Thus the finite polylogarithmic reduction itself, rather than only the
preceding integral transformations, has been checked independently.
For the total angular coefficient, taking
$z/Z=1-10^{-25}$ gives $d=8.8303688\times10^{-26}$.
Adding the analytically derived endpoint tail then yields
$\mathcal C_K(10)=10.058462232631976837745\ldots$.
Its difference from the original integral is $7.8\times10^{-48}$,
consistent with the $O(d^2\ln d)$ tail remainder. At $s=1.6$ the
interior endpoint check agrees to $2.8\times10^{-58}$.
These checks use only the amplitudes, soft integral, and algebraic
reductions derived here, with no published total cross section as input.

\subsection{Check of the inner-energy Laurent coefficients}
The script \nolinkurl{compton_real_uv_endpoint.py} verifies the
Laurent coefficients against the exact beta--hypergeometric
expressions at 80 decimal digits. This check evaluates analytic
functions and does not use numerical quadrature.

\subsection{Checks of the UV difference}
At $H=9$, the generated endpoint expression for this pole coefficient
contains six terms, with words drawn from
$G(r_\pm-1,-1;9)$, $G(r_\pm-1;9)$, and $G(-1;9)$, together with a
rational endpoint term. Evaluating this analytic expression gives
\begin{equation}
 \int_1^{10}\mathcal P_{\mathrm{UV}}(9,v)\,\dd v
 =27551.4550049022114364696640244483\ldots.
\end{equation}
The difference from independent quadrature is $4.4\times10^{-61}$ at
65-digit working precision. The quadrature uses the analytic local series
near $v=1$ to avoid subtracting individually divergent numerical terms.
The integrated result includes the complete difference with its Taylor
subtraction, rather than either divergent component in isolation.

The finite coefficient is reduced by the same rules to 26 terms.
Its origin expansions at two independent values of $H$ are
\begin{align}
 \mathcal F_{\mathrm{UV}}(9,1+\zeta)
 &=\left(\frac{24064}{729}-\frac{28672\ln3}{243}\right)\zeta+O(\zeta^2),\\
 \mathcal F_{\mathrm{UV}}(3/5,1+\zeta)
 &=\left(\frac{34816}{15}+\frac{114688}{135}\ln\frac53\right)\zeta+O(\zeta^2).
\end{align}
The integrated endpoint values are
\begin{align}
 \int_1^{10}\mathcal F_{\mathrm{UV}}(9,v)\,\dd v
 &=114586.9312643746457220611778438924\ldots,\\
 \int_1^{8/5}\mathcal F_{\mathrm{UV}}(3/5,v)\,\dd v
 &=504.4462882483088914154537366764\ldots.
\end{align}
At 65-digit precision, independent quadrature differs by
$6.3\times10^{-60}$ and $4.2\times10^{-62}$, respectively.
These values concern this angular difference sector before the overall
normalization and the other sectors are assembled.
The callable routine \texttt{uv\_difference\_closed} returns its pole or
finite coefficient from the generated polylogarithmic expression.
It performs no quadrature.

\subsection{Check of the UW pole integral}
At $H=9$ its integrand begins as $-1152t+O(t^2)$, and
\begin{equation}
 \int_0^1g_1(9,t)\,\dd t
 =-15.9227404613352952442617815523\ldots.
\end{equation}
At 60-digit working precision this expression agrees with independent
quadrature to $2.6\times10^{-56}$.

\subsection{Checks of the UW function basis}
The eight functions in the inner finite result have also been checked
against their real logarithm and dilogarithm definitions. At $H=9$ and
$t=0.1,0.5,0.9$, their polylogarithmic translations agree to better than
$2.4\times10^{-45}$ at 45-digit working precision. This includes the
difference of two dilogarithms with large negative arguments near $t=0$.
The identity is fixed by differentiation and its value at the origin,
so the numerical check does not choose a physical branch.

\subsection{Checks of the finite UW integral}
Its endpoint values before the common overall normalization are
\begin{align}
 \mathcal F_{UW}(9)
 &=-599.2785090120818758407150390268563\ldots,\\
 \mathcal F_{UW}(3/5)
 &=5.33372142619549103751931327587136\ldots.
\end{align}
Independent coordinate quadrature differs by $2.4\times10^{-74}$ and
$1.3\times10^{-65}$, respectively. These differences are consistent with
the estimated analytic tail remainder. The production routine
\texttt{real\_uw\_closed} returns the pole and finite coefficients using
only the polylogarithmic expression and convergent endpoint series.
It does not include the crossed angular sector.

\subsection{Check of the UW inner finite coefficient}
The independent numerical check integrates
the original $a$ expression only after subtracting its two endpoint
poles. It agrees with \eqref{eq:comp-real-UW-inner-closed} at an
interior point $H=0.6,t=0.3$ to better than $10^{-43}$ relatively.
The numerical integral is a check, not a definition of the answer.

\subsection{Checks of the soft-boundary conversion}
For the kernels $K=x^{-2}$ and $K=x^{-3}$ of
Section~\ref{sec:compton-soft-matching-check}, the comparison is as follows.
At $H=9$,
evaluation of the exact Mellin formulas at $\eta=\pm10^{-22}$ with
90-digit arithmetic reproduces these finite terms to more than 42 digits.
The exact cutoff formulas at $\zeta=10^{-30}$ agree with their displayed
asymptotic finite parts to more than 30 digits.  These evaluations test
analytic formulas and do not replace the integration by quadrature.

\subsection{Checks of the Bose subtraction correction}
The moments $P,A,B,C$ and finite parts $F_z,F_\eta$ use the definitions
of Section~\ref{sec:compton-soft-matching-check}. Here $H=s-1$,
and $D$ is the rational correction in
\eqref{eq:comp-real-bose-restoration}.
An independent evaluation at $H=9$ gives
\begin{align}
 P&=37497.9773085349211509285356121\ldots,\\
 A&=3813.01541537757944985908409943\ldots,\\
 B&=-1194.39517843100551019039928260\ldots,\\
 C&=875.197982980881880082536498852\ldots,\\
 F_\eta-F_z&=8697.24837987961462163725697550\ldots.
\end{align}
These values are evaluations of the closed moments, with no numerical
integration.  They provide the finite matching contribution of this
individual subtraction correction.

At $H=9$, the resulting closed endpoint value is
\begin{align}
 F_z[D]&=-34968.0992648172707249254370242076502947231\ldots,\\
 F_\eta[D]&=-26270.8508849376561032881800487118253811979\ldots.
\end{align}
Evaluations requested at 45 and 70 decimal digits agree in all the displayed
digits.  As an independent check only, quadrature of the one-dimensional
kernel up to $u=s-0.02$ agrees with its generalized-polylogarithm primitive
to an absolute difference below $1.1\times10^{-79}$ at 85-digit precision.
The finite endpoint value itself uses the analytic asymptotic expansion,
not quadrature or a fit to a cutoff sequence.

\subsection{Finite parts of the real-energy primitives}
Exact symbolic differentiation verifies the displayed GPL
primitives for all three sectors. As independent checks, at $H=9$
their $z$-cut finite parts are
\begin{align}
 F_{z,W}&=-199.1245079022571280440396316233659,\\
 F_{z,\ln s}&=371.2012216070278610540821023119374,\\
 F_{z,\mathrm{reg}}&=-13063.03114094500957535028266440329.
\end{align}
These numbers do not include the common factor $1/(8sH)$ or the
conversion from the $z$ cut to the symmetric analytic regulator.
Separate real integrations check the interior primitive boundary
differences to better than fifty decimal places at this point.
Continuing the full regular-sector GPL combination above and below
its artificial interior letter agrees to eighty decimal places.
These checks concern evaluation of the analytic primitives and do
not supply coefficients to them.

\subsection{Checks of the real soft residues}
The following checks evaluate the analytic residue moments and their
integrated sum. At $H=0.6,9,99$, separate 70-digit quadratures of these residue
moments agree with the analytic expressions to better than
$10^{-60}$ relatively.

At $H=9$ this correction is
\begin{equation}
 A_D=3813.01541537757944985908409942886783767351749\ldots.
\end{equation}
Evaluating the analytic expressions on both sides of
\eqref{eq:comp-real-integrated-pole-check} gives
\begin{center}
\begin{tabular}{rll}
\toprule
$H$ & Common value & Relative difference\\
\midrule
$0.6$ & $0.976756780072123721541847283885$ & $7.0\times10^{-68}$\\
$9$ & $730.423152821833269374946708743$ & $5.0\times10^{-70}$\\
$99$ & $50871.5845189161404834801869055$ & $6.0\times10^{-66}$\\
\bottomrule
\end{tabular}
\end{center}
This is a numerical check of the separately derived analytic
coefficients at three energies. It does not replace an identity
proof for arbitrary $H$, nor does it verify the finite total.

\subsection{Fixed-energy checks of endpoint subtraction}
At $H=9$ the pole integrand has the local expansion
\begin{equation}
 \mathcal P_{\mathrm{UV}}(9,1+\zeta)
 =-\frac{7168}{243}\zeta+O(\zeta^2).
\end{equation}
Thus its apparent origin poles cancel and the lower endpoint is regular.

As an independent check, the coefficient of its analytic-regulator pole
at $H=9$ integrates to
\begin{equation}
 P_{\mathrm{shift}}(9)
 =20\ln10+\frac{2913280}{59319}\ln\frac{11}{2}
 +\frac{5602208}{1521}.
\end{equation}
This term checks the subtraction bookkeeping independently of the
finite-angle polylogarithm evaluation.

\section{Checks of the analytic simplification}
These checks compare the simplified expressions with the original
independent sector results. The algebraic derivations above fix all
coefficients and boundary values before any numerical comparison.

\subsection{Checks of the classical root functions}
The reduction was also checked against the defining integrals at
$X=10^{-6}$, $10^{-2}$, $0.6$, $1$, $10$ and $10^4$, using 70-digit arithmetic.
The substitution $X=4\sinh^2(y/2)$ removes the square-root kernel
from those numerical checks. All 36 endpoint comparisons have absolute
differences below $1.2\times10^{-68}$. Independent differentiation
checks of all six classical expressions at five of these points also
reproduce their defining kernels. These numerical checks supplement
the primitive identities and boundary derivations above. They do not
supply any part of the analytic expressions.

\subsection{Checks of the simplified soft primitive}
At an interior endpoint, the old primitive equals
$\widetilde\Psi+4apf/(\beta E^2dk)$, which permits a direct check
without a numerical integral.  At $z=0.7Z$, evaluating both GPL
expressions with 65 decimal digits gives differences below
$2.6\times10^{-63}$ for $s=1.6$ and $s=10$.  These checks support
the implementation of the exact relation
\eqref{eq:comp-ck-pole-removal}.

\subsection{Checks of the collected soft kernel}
As a cross-check of this reorganization, numerical integration of the
regular kernel gives
\begin{align}
 \mathcal C_K(1.6)&=0.1712265509990548695712562109522182979661\ldots,\\
 \mathcal C_K(10)&=10.0584622326319768377445483326880109081643\ldots.
\end{align}
These values agree with the previous analytic soft expression.
The numerical integrations test the algebraic reorganization and do
not replace the analytic expression.

\subsection{Checks of the soft energy-variable integrals}
Direct quadrature of the defining
$I_1,I_2,I_7$ integrals at $X=10^{-3},0.6,9,100$ agrees with their
energy-variable analytic expressions to better than $3\times10^{-72}$
in absolute value for $I_2$ and $I_7$. The integral evaluations in this
check are not used to determine any coefficient or boundary constant.

\subsection{Checks of the integrated UW expression}
Evaluation of the fixed-alphabet result and the original finite GPL primitive gives
\begin{align} F_{UW}^{\eta}(0.6)&=5.33372142619549103751931327587135840926047\ldots,\\
F_{UW}^{\eta}(9)&=-599.278509012081875840715039026856308771543\ldots.\end{align}
The absolute differences are $1.4\times10^{-65}$ and $3.1\times10^{-63}$, respectively. These checks concern the $UW$ contribution. The conversion to \eqref{eq:comp-uw-integrated-compact} uses the exact coefficient and shuffle identities above.

\subsection{Checks of the matched real contribution}
Fresh evaluations of the original four GPL primitives at $s=1.6$ and $10$ agree with the endpoint formula before matching. Adding the original finite shifts gives agreement with \eqref{eq:comp-real-four-matched-compact} through the 45 significant digits recorded for those original values. These checks use our own integrated expressions. No published total is substituted.

\subsection{Checks of the four-sector endpoints}
The coefficients $A_4,D_4$ extracted from the collected primitive were
also checked against the Laurent expansions of the four original
integrands at $s=1.01$, $1.6$ and $10$, using 70-digit arithmetic.
This checks the endpoint signs and normalization independently of the
primitive coefficient collection.

\subsubsection{The conjugate-pair endpoint}
The function $\mathcal A_{\mathrm{pair}}(s)$ is defined in
\eqref{eq:comp-real-a-pair-endpoint}. Direct quadrature of $P'\ln B/P$ at six points from
$s=1.000001$ to $10^4$ agrees with this expression to an absolute
error below $6\times10^{-70}$ in 70-digit arithmetic.

\subsection{Checks of the direct-box basis conversion}
The independent direct-box conversion is derived in
Section~\ref{sec:comp-virtual-energy-basis}. As a check
of its implementation, evaluation at $s=1.6$ and $s=10$ differs from the
original closed endpoint representation by less than
$3.4\times10^{-78}$.  The exact derivative and boundary arguments,
rather than these two numbers, establish the basis transformation.


\section{Numerical evaluation of the published reference}
Only after completing the amplitude integration do we evaluate the
published total in Equation (6) of Lee, Schwartz and
Zhang~\cite{qed-processes-lsz}.  Its thirteen iterated-integral contributions
are combined into one integral in $t=\ln z$, so the lower-weight functions
at a common argument are computed once.  Near $z=0$, convergent power
series evaluate the lower-weight functions without cancellation between
constant terms.  This numerical integration evaluates the special
functions in the published reference.  It is not an input to the
independent derivation or to its closed endpoint formulas.

We compare evaluations with 75 and 90 decimal digits and lower endpoints
$t=-140$ and $t=-165$.  The omitted tails have analytic absolute bounds.
For example, for $0<z<1$ the lower-weight words obey
\begin{align}
 |G(-1,0;z)|&\leq z(1+|\ln z|),\\
 0\leq G(-2,-1;z)&\leq z^2/4,\\
 0\leq G(r,-1;z)&\leq z^{3/2}/3.
\end{align}
Here the letter $r$ denotes the kernel $1/\sqrt{z(z+4)}$, as in the
published formula, rather than the angular variable used earlier.
Multiplication by the outer kernels and integration of these elementary
bounds controls the lower-endpoint truncation.  The reported error
estimate adds these tail bounds, ten times the difference between the two
evaluations, and a rounding estimate based on working precision and the
absolute sizes of the individual coefficients.  It is an independent
numerical error estimate, not an interval-arithmetic proof and not an
estimate inferred from agreement with our result.
For the seven values $s=1.0002$, $1.02$, $1.6$, $10$, $100$, $2001$, and
$200001$, all reference error estimates are below $5.6\times10^{-58}$.
At $s=1.6$, evaluating each published iterated integral separately with
75-digit arithmetic and the same lower endpoint $t=-140$ differs from
the combined evaluation by $7.74\times10^{-75}$.  This checks the
combination independently of the comparison with our cross section.

\subsection{Evaluation of the reference words}
The implementation evolves the thirteen words in $\mathcal H$ and uses
the elementary expressions for $G(-1;X)$, $G(0,-1;X)$, $G(-1,0;X)$ and
$G(-1,-1;X)$ from Section~\ref{sec:comp-reference-elementary} for their
lower-weight derivatives. At
$\tau=\ln X$ the coefficient for a rational letter is
$X/(X-a)$ and that for $r$ is $\sqrt{X/(X+4)}$.
Each derivative involves a lower-weight word. All thirteen unknowns
vanish at $X=0$. Starting at $\tau=-36$ with zero initial values
neglects only endpoint powers, and repeating at $\tau=-40$ checks
this approximation. Large cancellations among the terms of $\mathcal H$
make the accuracy of the final correction weaker than the accuracy
of each individual $G$. The verification reports the error of the
combined cross section as well as that of the functions.


Numerically the double-precision implementation is intended for
$X\geq0.01$ and the tested energies, rather than arbitrarily close
to threshold or at arbitrarily large energy.

\section{Checks of the collected total}
At $X=0.6,9,99$, the collected formula gives
\begin{align}
C_1(0.6)&=0.39497522084295373429747165489072725825168951498178\ldots,\\
C_1(9)&=0.77515165401174580343102580969237768056470578000273\ldots,\\
C_1(99)&=0.35308838117224168652124506257313479270540259204215\ldots.
\end{align}
The absolute differences from the original independent analytic evaluations are $2.2\times10^{-79}$, $8.8\times10^{-81}$, and $1.1\times10^{-85}$, respectively. These are checks between two representations of our own calculation. No published total is used in the collection or in these comparisons.



\section{Integrated benchmarks}
The program \texttt{check\_compton\_nlo.py} tests four elementary
reductions of the $G$ functions at $X=0.07,0.5,1,10,20000$.
The maximum absolute error divided by $\max(1,|G|)$ is
$7.5\times10^{-14}$. Tightening both differential-equation
tolerances and lowering the starting point changes the correction
at the three tested PrimEx energies by less than $2.5\times10^{-8}$
of that correction. This is far below the experimental uncertainty.

For a head-on laser of energy $1.165\,\mathrm{eV}$, transform an
electron beam energy $E$ to the rest-frame photon energy by
\begin{equation}
 \omega=\omega_{\mathrm{laser}}
                  \frac{E+\sqrt{E^2-m^2}}m.
\end{equation}
Using $m=0.51099906\,\mathrm{MeV}$ and
$\alpha_{\mathrm T}^{-1}=137.0359895$ gives
Table~\ref{tab:compton-nlo-benchmark}. The reference values are
theoretical predictions from Table 3 of Denner and Dittmaier
\cite{qed-processes-denner}, used here to check the numerical
implementation. They are not experimental input to the PrimEx figure.
\begin{table}[!htbp]
 \centering
 \begin{tabular}{ccc}
 \toprule
 $E$ [GeV] & $100\sigma_1/\sigma_0$ [\%], evaluation
          & Reference~\cite{qed-processes-denner}, Table 3 [\%]\\
 \midrule
 4 & 0.0023092 & 0.0023\\
 6 & 0.0044867 & 0.0045\\
 8 & 0.0070869 & 0.0071\\
 \bottomrule
 \end{tabular}
 \caption{Inclusive photonic corrections with no acceptance cuts.
 The laser energy is $1.165\,\mathrm{eV}$ in all three rows.}
 \label{tab:compton-nlo-benchmark}
\end{table}

At center-of-mass energy
$\sqrt{s_{\mathrm{phys}}}=1\,\mathrm{GeV}$ the correction is $8.17\%$, which
reproduces the $8.2\%$ quoted in Reference~\cite{qed-processes-lsz}.
These comparisons verify the implementation of Equation (6) of
Lee, Schwartz and Zhang~\cite{qed-processes-lsz}.

\section{Experimental input and reproducibility}
The data source is Table 2 of the PrimEx publication
\cite{qed-processes-primex}, the same thin-carbon measurements used
in the book's tree-level comparison. An unchanged copy of the raw table is included as
\nolinkurl{data/primex_carbon_raw.dat}. Its columns give the run period,
photon energy in GeV, cross section in mb per electron, and separate
statistical and systematic uncertainties in the same cross-section
units. The 29 values are copied without fitting or rescaling.

The comparison uses $m=0.51099895\,\mathrm{MeV}$,
$\alpha_{\mathrm T}^{-1}=137.035999084$, and
$1\,\mathrm{GeV}^{-2}=0.3893793721\,\mathrm{mb}$.
The plotted error is the quadrature sum of the two tabulated errors.
The published cross sections include background, luminosity and
acceptance corrections and are not Born-unfolded quantities.
A detector simulation is therefore not reapplied to the reported
total. The experiment's extraction already uses a radiative-event
model. The comparison does not redo that extraction.

At $5\,\mathrm{GeV}$ the tree and inclusive NLO predictions are
$0.264648\,\mathrm{mb}$ and $0.273692\,\mathrm{mb}$ respectively.


All calculation programs are in \texttt{scripts/}.
\texttt{compton\_nlo.py} evaluates the exact total in Equation (6)
of Lee, Schwartz and Zhang~\cite{qed-processes-lsz}.

\texttt{generate\_compton\_nlo\_comparison.py} evaluates the independently
integrated total through
\nolinkurl{compton_total_independent_closed.py}, without calling the
published formula. It writes sampled theory,
experimental points and their errors as \texttt{.dat} files in the
local \texttt{tex/figures/} directory. The vector figure is drawn by PGFPlots
during the chapter build. Verification programs write their numeric
reports in the \texttt{reports/} directory. The build instructions give
the commands and required Python libraries.

Figure~\ref{fig:compton-nlo-primex} compares the measured totals with
both predictions. The inclusive radiative correction raises the curve
over this energy range. The lower panel displays this change relative
to the Born prediction using the same normalization for data and theory.
\providecommand{\comptonfiguredir}{figures}
\begin{figure}[!htbp]
 \centering
 \begin{tikzpicture}
  \begin{axis}[comptoncomparison,name=total,
   ylabel={$\sigma$ [mb/electron]},xticklabels={}]
   \addplot[comptontree] table[x=energy,y=born]
    {\comptonfiguredir/compton_nlo_total.dat};
   \addlegendentry{Klein--Nishina tree}
   \addplot[comptonnlo] table[x=energy,y=nlo]
    {\comptonfiguredir/compton_nlo_total.dat};
   \addlegendentry{NLO: virtual + real}
   \addplot[comptonfirst,
    error bars/.cd,y dir=both,y explicit]
    table[x=energy,y=sigma,y error=error]
    {\comptonfiguredir/compton_nlo_primex_1.dat};
   \addlegendentry{PrimEx, Run 1}
   \addplot[comptonsecond,
    error bars/.cd,y dir=both,y explicit]
    table[x=energy,y=sigma,y error=error]
    {\comptonfiguredir/compton_nlo_primex_2.dat};
   \addlegendentry{PrimEx, Run 2}
  \end{axis}
  \begin{axis}[comptoncomparison,
   at={(total.south west)},anchor=north west,yshift=-0.25cm,
   height=4.2cm,xlabel={$\omega$ [GeV]},ylabel={Ratio to tree},
   ymin=0.975,ymax=1.08]
   \addplot[comptontree,domain=4.35:5.55]{1};
   \addplot[comptonnlo] table[x=energy,y expr=1+\thisrow{percent}/100]
    {\comptonfiguredir/compton_nlo_total.dat};
   \addplot[comptonfirst,
    error bars/.cd,y dir=both,y explicit]
    table[x=energy,y=ratio_born,y error=error_born]
    {\comptonfiguredir/compton_nlo_primex_1.dat};
   \addplot[comptonsecond,
    error bars/.cd,y dir=both,y explicit]
    table[x=energy,y=ratio_born,y error=error_born]
    {\comptonfiguredir/compton_nlo_primex_2.dat};
  \end{axis}
 \end{tikzpicture}
 \caption{The PrimEx thin-carbon cross sections per electron,
 compared with the tree prediction and our analytically integrated
 NLO result, including virtual corrections and double Compton emission.
 Points are from Table 2 of Reference~\cite{qed-processes-primex}.
 Run 1 and Run 2 denote the two data-taking periods.
 Error bars combine
 the reported statistical and systematic uncertainties in quadrature.
 The lower panel divides both data and theory by the tree prediction.
 }
 \label{fig:compton-nlo-primex}
\end{figure}


\section{Independent numerical check of the real contribution}
The numerical check integrates the original six-diagram spin sum,
using explicit $4\times4$ Dirac matrices and transverse photon
polarizations. It does not evaluate the GPL primitives or insert
the published total cross section. We use the ordered-photon,
locally subtracted integral in
\eqref{eq:comp-supp-real-subtracted}. Thus its output is precisely
$\mathcal C_{\mathrm R}$ in
\eqref{eq:comp-supp-independent-total}, including the
$\mathcal A_{\mathrm{loc}}\ln\omega_*$ term. No finite photon-energy cut remains.

To resolve the two electron directions, the angular sampling uses
a normalized mixture of two positive densities. With
$\beta=q/E$ and $\widehat{\bm v}_{i,f}$ the incoming and elastic
outgoing electron directions, define
\begin{equation}
 g_{i,f}(\widehat{\bm n})=
 \frac{\beta}{2\pi\ln s}
 \frac1{1-\beta\widehat{\bm v}_{i,f}\cdot\widehat{\bm n}},
 \qquad g=\tfrac12(g_i+g_f).
\end{equation}
Each density integrates to one because
$\ln[(1+\beta)/(1-\beta)]=\ln s$. Each channel is mapped to a
unit square by using the logarithm of
$1-\beta\widehat{\bm v}_{i,f}\cdot\widehat{\bm n}$ and a uniform
azimuth. The two channel contributions are averaged with weight
$1/g$. The elastic polar angle uses the analogous logarithmic map
for $E+qc_2$, and the last coordinate is $\omega/\omega_*$.
The resulting finite four-dimensional integral is evaluated by
adaptive cubature.

At each energy we repeat the calculation with increasingly small
requested tolerances. The numerical uncertainty assigned to the
last result is twice the largest of the last two cubature error
estimates and the change between those two results. It is an
empirical convergence estimate, not a rigorous bound. A run that
exhausts its evaluation budget is explicitly distinguished from a
run satisfying the requested tolerance. For the low-energy endpoints and $s=2001$, the matrix and
kinematic arithmetic use \texttt{long double}, while the cubature
weights and accumulated results use double precision. At $s=200001$
the double-precision calculation is also repeated independently with
\texttt{long double} at the same requested tolerance.
The saved audit records identify the source, compiler command,
binary hash, requested tolerances, evaluation counts, and every
refinement result.

For comparison with the symmetric analytic regulator, if its real
coefficient is $A/\eta+F$, the precise conversion is
\begin{equation}
 \mathcal C_{\mathrm R}=F-A\ln[2\sqrt{s}(s-1)].
\end{equation}
This follows from $xy\to2\sqrt{s}(s-1)\omega$ at a soft endpoint.
The factor for the identical photons is already accounted for by
the ordering. Consequently this numerical result can be added
directly to the independently integrated
$\mathcal C_{\mathrm{VS}}$, without importing an analytic finite
constant from the total cross section.

At the highest energy the successive real coefficients are
\begin{equation}
 0.01360181858,\quad 0.01361650195,\quad
 0.01362436493,\quad 0.01362842773.
\end{equation}
The refinement changes decrease from $1.47\times10^{-5}$ to
$7.86\times10^{-6}$ and then $4.06\times10^{-6}$. Accordingly the
assigned uncertainty is $8.13\times10^{-6}$, larger than the final
cubature error estimate. At the penultimate tolerance, the independent
double-precision and \texttt{long double} runs each use
$4\,157\,469$ evaluations and differ by $2.31\times10^{-15}$.
The nodes were not recorded, so equal evaluation counts do not
establish an identical adaptive grid. This precision comparison is
consistent with integration convergence dominating the reported
uncertainty. These numerical tolerances do not limit the precision
of the independently evaluated analytic formula.


\section{Three-way comparison over the energy range}
\label{sec:comp-independent-total-comparison}
The coefficient compared here is
\begin{equation}
 C(s)=C_1(s-1)=\frac{m^2\sigma_1}{\alpha_{\mathrm T}^3},\qquad
 s=\frac{s_{\mathrm{phys}}}{m^2}=1+\frac{2\omega}{m}.
\end{equation}
The photon energy $\omega$ is measured in the initial electron rest
frame. The seven energies span $10^{-4}\leq\omega/m\leq10^5$.
The first value at each energy in
Table~\ref{tab:comp-independent-total-50} is our independent analytic
coefficient, freshly evaluated from the collected expression in
Equation~\eqref{eq:comp-total-energy-collected}. Its 50 displayed digits
also agree with the original assembled sector primitives.
All 50 displayed significant digits agree between two runs at
65 and 85 requested decimal digits, each with 40 guard digits.
This agreement checks numerical stability. The endpoint series used
to evaluate our primitives have truncation-error estimates, but a
rigorous bound on every omitted tail has not been established.
The displayed digits therefore do not constitute a certified
50-digit error bound.

The two subsequent entries are absolute coefficient differences,
\begin{equation}
 \Delta C_{\mathrm{num}}=C_{\mathrm{num}}-C(s),\qquad
 \Delta C_{\mathrm{ref}}=C_{\mathrm{ref}}-C(s).
\end{equation}
The differences use the full saved analytic value before rounding it
to the 50 digits displayed in the table.
For $C_{\mathrm{num}}$, the virtual and soft terms are integrated
numerically over the elastic angle from the pointwise amplitude.
The real term is a separate four-dimensional phase-space integral
of the six-diagram double-Compton matrix element, with its local
soft subtraction. Neither calculation calls the analytic total
or its integrated GPL primitives.

The real quadrature error estimate includes both the integrator's
reported error and the change upon refinement. The virtual estimate
includes angular refinement and an arithmetic accuracy floor.
We add these estimates rather than treating them as independent
statistical errors. Thus the $\pm$ entries are numerical error
estimates, not confidence intervals or experimental uncertainties.
The numerical integration is not required to reproduce 50 digits.

The reference coefficient is Equation (6) of Lee, Schwartz and
Zhang~\cite{qed-processes-lsz}. It is evaluated separately, after
our analytic result is fixed. Its error estimate uses two working
precisions, changes in the special-function integration endpoint,
an analytic bound on the omitted tail, and a rounding scale based
on the sizes of its terms. It does not use the difference from our
answer to estimate its own error. This comparison tests numerical evaluation of the two analytic
expressions. Section~\ref{sec:compton-analytic-equivalence} separately
establishes their exact equality by comparing every coefficient and
fixing their common physical continuation.

\begin{table}[!htbp]
\centering
\comptonprecisiontable
\begin{tabular}{rll}
\toprule
$\omega/m$ & Calculation & Coefficient or difference\\
\midrule
$0.000100$ & Our analytic & $0.00000062423288400014031134169341516503703552376453828153$ \\
 & Numerical difference & $5.30\times10^{-15}\;\pm\;2.0\times10^{-12}$ \\
 & Reference difference & $-4.03\times10^{-80}\;\pm\;5.6\times10^{-58}$ \\
\addlinespace[3pt]
$0.0100$ & Our analytic & $0.0029187564665568768995602481244244431628302067748237$ \\
 & Numerical difference & $1.11\times10^{-10}\;\pm\;2.0\times10^{-8}$ \\
 & Reference difference & $2.37\times10^{-86}\;\pm\;6.4\times10^{-65}$ \\
\addlinespace[3pt]
$0.300$ & Our analytic & $0.39497522084295373429747165489072725825168951498178$ \\
 & Numerical difference & $2.72\times10^{-9}\;\pm\;6.8\times10^{-7}$ \\
 & Reference difference & $8.37\times10^{-90}\;\pm\;1.1\times10^{-69}$ \\
\addlinespace[3pt]
$4.50$ & Our analytic & $0.77515165401174580343102580969237768056470578000273$ \\
 & Numerical difference & $5.71\times10^{-8}\;\pm\;1.2\times10^{-6}$ \\
 & Reference difference & $2.40\times10^{-91}\;\pm\;3.7\times10^{-72}$ \\
\addlinespace[3pt]
$49.5$ & Our analytic & $0.35308838117224168652124506257313479270540259204215$ \\
 & Numerical difference & $-3.86\times10^{-8}\;\pm\;1.2\times10^{-6}$ \\
 & Reference difference & $-6.91\times10^{-89}\;\pm\;1.0\times10^{-71}$ \\
\addlinespace[3pt]
$1.00\times10^{3}$ & Our analytic & $0.070463896610137198990117049003399503906848383360028$ \\
 & Numerical difference & $3.33\times10^{-6}\;\pm\;1.5\times10^{-5}$ \\
 & Reference difference & $-4.10\times10^{-87}\;\pm\;6.7\times10^{-71}$ \\
\addlinespace[3pt]
$1.00\times10^{5}$ & Our analytic & $0.0028704528832006358222188994983077446681129388945786$ \\
 & Numerical difference & $-1.34\times10^{-6}\;\pm\;8.1\times10^{-6}$ \\
 & Reference difference & $-9.38\times10^{-87}\;\pm\;5.5\times10^{-70}$ \\
\bottomrule
\end{tabular}
\caption{Three independent evaluations of the NLO Compton correction
$\sigma_1=(\alpha_{\mathrm T}^3/m^2)C_1$, from low to high incident
energy in the electron rest frame. Our analytic coefficient $C_1$ is
evaluated from the collected formula and printed to 50 significant digits. Each difference subtracts this
coefficient from the indicated calculation. The $\pm$ entries estimate
numerical error, not experimental uncertainty. The reference is
Equation (6) of Lee, Schwartz and Zhang~\cite{qed-processes-lsz}.}
\label{tab:comp-independent-total-50}
\end{table}


